\section{Architecture}

\begin{figure*}[t]
\centering
\includegraphics[width=0.95\textwidth]{fig/platform_architecture.pdf}
\caption{Malak Platform architecture showing the complete pipeline from training to deployment. The framework integrates Python-based training and compression tools, multi-compiler support, EdgeIR optimization passes, and a portable C++ runtime with hardware abstraction.}
\label{fig:architecture}
\end{figure*}

Figure \ref{fig:architecture} illustrates the Malak Platform architecture, organized into seven major components that span the complete edge AI lifecycle.

\subsection{Core Components}

\subsubsection{Model Cards and Schemas}
Model cards provide YAML-based metadata capturing task specifications, model sizes, latency/energy budgets, and calibration ranges. Protocol Buffer schemas define language-agnostic interfaces for features, telemetry, and on-device events, ensuring interoperability across Python tools and C++ runtime.

\subsubsection{Privacy and Policy Framework}
JSON-based policies encode privacy requirements (data retention, anonymization), fallback strategies (cloud escalation triggers), and compliance rules. This declarative approach separates policy from implementation, enabling auditable deployment in regulated domains.

\subsection{Python Training and Compression Pipeline}

\subsubsection{Training Module}
Built on PyTorch, our training pipeline incorporates:
\begin{itemize}
\item \textbf{Knowledge distillation}: Student models learn from teacher outputs, reducing parameters while preserving accuracy
\item \textbf{Quantization-aware training}: Models trained with fake quantization to adapt to INT8/INT4 precision
\item \textbf{Data augmentation}: Domain-specific transforms for robust generalization
\end{itemize}

\subsubsection{Compression Toolkit}
Post-training optimizations include:
\begin{itemize}
\item \textbf{Symmetric/asymmetric quantization}: Per-tensor or per-channel with calibration-based range selection
\item \textbf{Structured pruning}: Channel-level sparsity for hardware efficiency
\item \textbf{Export formats}: ONNX and TFLite for broad compatibility
\end{itemize}

\subsubsection{Multi-Compiler Support}
We integrate three compilation backends:
\begin{itemize}
\item \textbf{TVM}: AutoTVM-based autotuning for embedded ARM/RISC-V targets
\item \textbf{MLIR}: Custom dialects for fine-grained control over quantization and layout
\item \textbf{XLA}: TPU/Edge TPU acceleration paths
\end{itemize}

\subsection{EdgeIR: Intermediate Representation}

EdgeIR serves as the platform's common format, designed for:
\begin{itemize}
\item \textbf{Compactness}: FlatBuffers serialization minimizes binary size
\item \textbf{Optimization}: Graph-level passes (fusion, constant folding, layout transforms)
\item \textbf{Portability}: Schema-driven code generation for C++ runtime
\end{itemize}

Optimization passes include:
\begin{enumerate}
\item Batch normalization and ReLU fusion
\item Constant folding and dead code elimination
\item Layout transformation (NHWC $\leftrightarrow$ NCHW)
\item Quantization parameter propagation
\end{enumerate}

\subsection{C++ Runtime and Kernels}

\subsubsection{Graph Executor}
The static graph executor performs:
\begin{itemize}
\item Single-pass memory planning with in-place operations
\item Operator scheduling optimized for cache locality
\item Zero-copy inference pipeline
\end{itemize}

\subsubsection{Optimized Kernels}
Hand-tuned INT8/INT4 kernels for:
\begin{itemize}
\item 2D convolution (im2col + GEMM, Winograd)
\item Depthwise separable convolution
\item Matrix multiplication (blocking, vectorization)
\item Activations (ReLU, ReLU6, hardswish) with LUT acceleration
\item Pooling (max, average)
\end{itemize}

ARM NEON intrinsics provide 4-8� speedup over naive C implementations.

\subsubsection{Hardware Abstraction Layer}
Backend plugins support:
\begin{itemize}
\item \textbf{CPU}: ARM Cortex-M/A with NEON/Helium extensions
\item \textbf{CUDA}: NVIDIA Jetson edge devices
\item \textbf{NPU}: Dedicated neural accelerators (e.g., Ethos-U55)
\item \textbf{DSP}: Hexagon, CEVA architectures
\end{itemize}

\subsection{I/O and Telemetry}

\subsubsection{Sensor Integration}
Zero-copy DMA pipelines stream camera frames and sensor data directly to inference buffers, eliminating unnecessary memory copies.

\subsubsection{Telemetry System}
On-device telemetry captures:
\begin{itemize}
\item Per-layer latency and memory consumption
\item Energy measurements (via on-chip PMU or external monitors)
\item Input statistics for drift detection
\end{itemize}

\subsection{Monitoring and Production Tools}

\subsubsection{Drift Detection}
Statistical tests (KL divergence, population stability index) detect distribution shifts requiring model retraining.

\subsubsection{Accuracy Canaries}
Periodic ground-truth validation with flagged samples ensures model degradation alerts.

\subsubsection{CLI Interface}
Unified \texttt{edgeai} command-line tool orchestrates:
\begin{verbatim}
edgeai build --model jellyfish_v1.yaml
edgeai quantize --bits 8 --calib-data ./data
edgeai flash --target stm32h7
edgeai profile --device cortex-m7
\end{verbatim}

\subsection{Mojo Performance Layer}

For critical hot paths, Mojo kernels provide:
\begin{itemize}
\item MLIR-compatible quantized matmul microkernels
\item SIMD-friendly memory layouts
\item Compile-time specialization for target ISAs
\end{itemize}

This multi-language approach balances Python's expressiveness, C++'s portability, and Mojo's peak performance.
